\section{Evaluation}
\label{sec:evaluation}

This section asks three questions of \textsc{IntC}, in the order of the
subsections that answer them.
\begin{enumerate}
  \item Does \textsc{IntC} work on code extracted from real-world
  projects?
  \item Do as-if compilation of \(P_1\), and the oracle, budget filter, and
  \textsc{TheoryLower} score inside \textsc{IntentCompile}, behave as
  claimed?
  \item Under the same lowering backend that compiled \(P_1\), does
  \textsc{IntC} emit a candidate that is cheaper on
  \(\mathcal{O}\) than \(P_1\), or the same program?
\end{enumerate}
A last subsection names misses of \textsc{IntC}.
Runs use the machine and compilers of
Section~\ref{sec:implementation}. Tables are emitted by
\texttt{python reproduce.py} in the artifact
(\S\ref{sec:data-availability}). Solver times vary across runs. The
admit, timeout, and counterexample outcomes, and the rewrite names, are the
claims.

\subsection{Extracts from Real-World Projects}
\label{sec:catalog_table}

Each row of Table~\ref{tab:catalog} is one run of \textsc{IntC}
on a contract that states slack a developer could grant. How each extract
was cut from the upstream tree is in the artifact.

\begin{table}
\caption{One \textsc{IntC} run per row. Columns are the extract,
the catalog rule that fired, the harness bound \(K\), CBMC on \(P_1\) and
on the admitted program, and how that program compares to \(P_1\).}
\label{tab:catalog}
\begin{center}
\scriptsize
\begin{tabular}{llllll}
\toprule
Case & Rewrite & $K$ & $V$ on $P_1$ & $V$ on winner & As-if \\
\midrule
DPDK ring & pow2-mask & 4 & unknown/120.0s & proved/2.0s & not as-if \\
DPDK pad & pow2-mask & 4 & unknown/120.0s & proved/0.6s & unexposed \\
ClickHouse & pow2-mask & 4 & unknown/180.2s & proved/63.7s & not as-if \\
Zephyr & pow2-shift & 1 & proved/0.5s & proved/0.7s & not as-if \\
gain & pow2-mul & 1 & proved/0.1s & proved/0.1s & not as-if \\
PARSEC & aos-to-soa & 6 & proved/0.6s & proved/0.6s & layout \\
Flecs & predicated-update & 10 & proved/0.1s & proved/0.1s & layout \\
MachSuite & halo-1d & 8 & proved/0.1s & proved/0.1s & layout \\
CMSIS-DSP & align-loop-bound:BLOCK & 12 & proved/3.3s & proved/7.0s & not as-if \\
OGRE & narrow-int-field & 6 & proved/0.1s & proved/0.1s & layout \\
Box2D & squared-compare & 1 & proved/13.7s & proved/0.2s & not as-if \\
Opus & narrow-accum & 6 & proved/0.7s & proved/12.0s & not as-if \\
OpenCV & luma-q15 & 6 & unknown/60.0s & proved/0.1s & not as-if \\
PARSEC fat & aos-to-soa & 6 & proved/0.6s & proved/0.7s & layout \\
\bottomrule
\end{tabular}
\end{center}

\end{table}

\textbf{Case.} The cell names a tagged upstream situation whose hot loop we
extracted, except gain, which is a multiply-to-shift written for this
catalog. The DPDK ring and pad baselines are the 1000-slot sketches of
Section~\ref{sec:overview_intent}, not the library sources. ClickHouse's
baseline uses 10007 buckets, a prime, not the upstream power of two.
PARSEC and PARSEC fat are the same Euler step on four particles, the
latter on a 64-byte record that includes four explicit pad fields.
CMSIS-DSP is a two-tap FIR.

\textbf{Rewrite.} The cell is the catalog rule that produced the admitted
program. \textsc{TheoryLower} contributes remainder-to-mask, division-to-shift,
and multiply-to-shift. \textsc{GeneratePool} contributes the representation
rewrites and Rec.~601. In table order:
\label{sec:gold_rewrites}
\begin{itemize}
  \item DPDK ring: remainder-to-mask. CEGIS selects 1024.
  \textsc{GeneratePool} also admits pad-and-mask into that pool at \(K=4\).
  \item DPDK pad: remainder-to-mask on \texttt{STORAGE}, with
  \texttt{CAPACITY} still 1000. Fullness uses
  \(\texttt{tail}-\texttt{head}\) against \texttt{CAPACITY}, so enqueue
  fails at 1000 live packets. The extra 24 cells are a budget cost.
  \textsc{TheoryLower} and CEGIS produce that mask on the already-named
  \texttt{STORAGE}.
  \label{sec:unexposed_eval}
  \item ClickHouse: remainder-to-mask. \(\mathcal{I}\) requires at least
  10007 buckets. CEGIS selects 16384, the admissible power of two nearest
  that floor.
  \item Zephyr: \(\texttt{tick\_us / DIVISOR}\) to a shift. Relative error
  at most \(2.5\%\) versus \(\texttt{tick}/1000\) on 32-bit ticks at least
  640000, because integer division by 1024 overshoots that bound near
  600000. That floor is part of \(\mathit{Pre}\). CEGIS selects
  \(\texttt{DIVISOR}\) 1024 and \(\texttt{SHIFT}\) 10.
  \label{sec:decimation_eval}
  \item Gain: \(\texttt{sample * GAIN}\) with \(\texttt{GAIN} \in [8,16]\)
  to a left shift by 3. The harness asserts equality with
  \(\texttt{sample * GAIN}\) on 16-bit samples, so the admit is not a
  multiply by the baseline's 10.
  \label{sec:scale_eval}
  \item PARSEC: array-of-structs to structure-of-arrays.
  \label{sec:soa_eval}
  \item Flecs: predicated update on a 0/1 live flag.
  \item MachSuite: 1-D halo pad.
  \item CMSIS-DSP: loop bound aligned to a multiple of 4.
  \(\mathcal{I}\) is \(\texttt{BLOCK} \ge 5\) with search box
  \(\texttt{BLOCK} \le 64\); CEGIS selects 8.
  \item OGRE: 32-bit space tag narrowed to \texttt{uint8\_t}.
  \item Box2D: \(\texttt{sqrtf(d2) < R}\) to \(\texttt{d2 < R*R}\).
  \(\mathit{Post}\) is the integer disk test \(dx^2+dy^2 < R^2\); the
  \(\texttt{sqrtf}\) in \(P_1\) is not that test. The timed bench is
  \(10^6\) float pairs, outside \(\models_K\).
  \item Opus: \(\texttt{double}\) accumulator narrowed to \texttt{float}.
  The harness uses \(\mathit{LEN}=4\) operands in \(\{0,1,2,3\}\). The
  timed bench is \(L=256\) fractional floats, outside \(\models_K\).
  \item OpenCV: Rec.~601 float coefficients to a Q15 shift. The harness
  is that Q15 formula within one gray level, not the float Rec.~601 sum.
  \(\mathit{NPX}=4\) in the harness; the timed bench is \(1920\times 1080\),
  outside \(\models_K\).
  \item PARSEC fat: array-of-structs to structure-of-arrays.
\end{itemize}

{\bfseries\boldmath\(K\).} The cell is the harness bound for \(\models_K\): how many
operations or loop unwindings CBMC covers on that row
(\S\ref{sec:verifier}). Ring, pad, and ClickHouse use four operations from
an empty container. Zephyr, gain, and Box2D use one conversion, sample, or
compare. The rest are loop unwindings of the extract: 6 on PARSEC's four particles
and on OGRE, Opus, OpenCV, and PARSEC fat; 8 on MachSuite; 10 on Flecs;
12 on the two-tap FIR.

{\bfseries\boldmath\(V\) on \(P_1\) and on the winner.} Those cells are that CBMC
query. \(\mathsf{unknown}\) is a timeout. PARSEC and PARSEC fat prove the
Euler step at \(n=4\). CMSIS-DSP proves the two-tap filter. When both
programs prove, a longer time on the rewrite is the cost of that query.
Table~\ref{tab:catalog}'s Opus and CMSIS-DSP rows are the two where both
prove and the rewrite is the slower query.

\textbf{As-if.} The cell classifies the admitted program against \(P_1\).
\emph{not as-if} means a named parameter or a precision choice moved.
\emph{unexposed} means the backing store grew and caller-facing capacity
did not. \emph{layout} means a representation rewrite the harness treats as
unobservable.

\subsection{Behavior as Claimed}
\label{sec:microbench}

This subsection answers as-if compilation of \(P_1\), then the oracle, the
budget filter, and the \textsc{TheoryLower} score, in that order.

\textbf{As-if compilation of \(P_1\).}
The DPDK ring maps a sequence number \(i\) to a slot by remainder, and in
\(P_1\) that remainder is \(i \bmod 1000\). MSVC~\texttt{/O2} still emits a
reciprocal multiply for \texttt{i \% 1000} and \texttt{and ecx, 1023} for
\texttt{i \% 1024}, and it does not change 1000 to 1024. A microbenchmark of
that slot expression alone uses 15 interleaved trials of
\(2\cdot 10^8\) iterations and two seeds. A latency-bound chain measures
2.67\,ns versus 0.39\,ns, a factor of 6.8, and independent operations
measure 0.52\,ns versus 0.17\,ns, a factor of 3.1. That upper bound has no
queue memory traffic.

\label{sec:ablations}
Table~\ref{tab:ablations} is one switch at a time on that ring's contract
for the oracle and the budget.

\begin{table}
\caption{Four switches on the DPDK ring: FIFO oracle, \(\models_K\),
occupancy-only \(\mathit{Post}\), and an 8016-byte ceiling.}
\label{tab:ablations}
\begin{center}
\scriptsize
\begin{tabular}{p{0.34\linewidth}p{0.58\linewidth}}
\toprule
Switch & Outcome \\
\midrule
Contract oracle (FIFO harness) & proved/0.6s \\
Admission is $\models_K$ & K=4 proved/0.7s; K=10 counterexample/17.3s \\
Occupancy-only vs FIFO $\mathit{Post}$ & capacity-only proved/0.5s; FIFO counterexample/0.4s \\
BudgetNormalize at 8016\,B & emit P1 \\
\bottomrule
\end{tabular}
\end{center}

\end{table}

\textbf{Oracle.} The FIFO harness proves the admitted mask in 0.6\,s.
Admission is \(\models_K\): the colliding mask \(\texttt{\& 999}\) is
proved at \(K=4\) and refuted at \(K=10\), because slots collide once nine
or more elements are live.
\label{sec:bounded_gap}
An occupancy-only \(\mathit{Post}\) proves a LIFO stack that the FIFO
harness refutes.

\textbf{Budget filter.}
\label{sec:budget_eval}
\textsc{BudgetNormalize} at 8016 bytes, the 1000-slot baseline: CEGIS still
proves the 1024-slot mask at 8208 bytes. Algorithm~\ref{alg:budgetnormalize}
drops it and the emit is \(P_1\). The 1000-slot remainder still times out
at that \(K\).

\textbf{\textsc{TheoryLower} score.}
\label{sec:score_eval}
Forty remainder-site variants, on a ring at \(K=4\) and a hash set at
\(K=3\), have Spearman's \(\rho = 0.11\) between that score
(\S\ref{sec:implementation}) and CBMC time at a 120\,s timeout.
Every timeout still has a non-power-of-two remainder in the program, and
proved times are flat, so a score at or below the threshold is not a
decision that \(V\) will finish. A score that predicted CBMC time is a
different score (\S\ref{sec:discussion}). When TheoryLower finds no
remainder or division hole, it plants no seed: PARSEC's Euler step is that
case, and \textsc{GeneratePool} still admits structure-of-arrays from
\(P_1\).

\subsection{Cheaper on \(\mathcal{O}\), or \(P_1\)}
\label{sec:ring_eval}

This subsection answers whether ranking under the same lowering backend
emits a candidate cheaper on \(\mathcal{O}\) than \(P_1\), or \(P_1\).
Table~\ref{tab:timed} is the seven contracts that name a \texttt{bench.c}.
Min, max, and quartiles are in the artifact JSON.

\begin{table}
\caption{Contracts that name a \texttt{bench.c}, ordered by emit ratio,
\(P_1\) last. Columns are the extract, \(K\) on the timed harness, CBMC on
the admitted program, median wall-clock ns per unit of work of \(P_1\) and
of the admit under MSVC~\texttt{/O2}, and emit.}
\label{tab:timed}
\begin{center}
\scriptsize
\begin{tabular}{llllll}
\toprule
Case & $K$ & $V$ on admitted & $P_1$ ns & admitted ns & emit \\
\midrule
Box2D & 1 & proved/0.3s & $0.702 \pm 0.024$ & $0.108 \pm 0.015$ & $6.5\times$ \\
PARSEC fat & 6 & proved/0.6s & $1.203 \pm 0.039$ & $0.346 \pm 0.012$ & $3.5\times$ \\
ClickHouse & 4 & proved/62.4s & $9.455 \pm 0.097$ & $4.000 \pm 0.043$ & $2.4\times$ \\
DPDK ring & 10 & proved/47.8s & $1.128 \pm 0.017$ & $0.487 \pm 0.008$ & $2.3\times$ \\
Opus & 6 & proved/11.9s & $0.440 \pm 0.005$ & $0.273 \pm 0.003$ & $1.6\times$ \\
OpenCV & 6 & proved/0.1s & $0.840 \pm 0.011$ & $0.560 \pm 0.014$ & $1.5\times$ \\
CMSIS-DSP & 12 & proved/6.7s & $0.460 \pm 0.016$ & $0.927 \pm 0.062$ & $P_1$ \\
\bottomrule
\end{tabular}
\end{center}

\end{table}

\textbf{Case.} The cell is the Table~\ref{tab:catalog} extract whose timed
contract names that driver.

{\bfseries\boldmath\(K\).} The cell is the timed harness bound. Six of the
seven match Table~\ref{tab:catalog}. The DPDK ring is \(K=10\) here and
\(K=4\) in that table. Table~\ref{tab:catalog} used four operations from
an empty queue. This row uses ten, the bound that refutes
\texttt{\& 999} in Section~\ref{sec:bounded_gap}.

{\bfseries\boldmath\(V\) on admitted.} The cell is CBMC of the admitted
rewrite on that timed harness, including CMSIS-DSP where emit is \(P_1\)
and the cell is still the slower \(\texttt{BLOCK}=8\) admit. The ring row
is that query at \(K=10\). Table~\ref{tab:catalog}'s ring cell is the same
rewrite at \(K=4\). Solver seconds move across runs; the proved/unknown
outcome is the claim.

{\bfseries\boldmath\(P_1\) ns and admitted ns.} Those two cells are
\(\mathcal{O}\): wall-clock nanoseconds per unit of work on that
\texttt{bench.c} (pair, update, lookup, enqueue, mac, pixel, or FIR
block), lower is better. The stage compiles both programs with
MSVC~\texttt{/O2} and takes the median of 31 trials after warmup. A later
candidate replaces the current best only if it is faster by more than its
sample standard deviation.

\textbf{Emit.} The cell is the ratio of those medians, or \(P_1\) when the
admit is not faster by more than its sample standard deviation.
In table order:
\begin{itemize}
  \item Box2D: in-range pairs. MSVC~\texttt{/O2} still emits
  \(\texttt{sqrtss}\) from the baseline, so the squared compare is the
  ranking win. The listing is in the artifact.
  \item PARSEC fat: the Euler step at \(N=100000\). The catalog query
  still proves \(n=4\).
  \item ClickHouse: fill 4000 keys, then \texttt{contains}. The ratio
  mixes a larger table and a lower load factor with the mask.
  \label{sec:hash_eval}
  \item DPDK ring: a dependent enqueue/dequeue chain through a partly
  full queue. The two programs disagree on the 1001st consecutive
  enqueue. The isolated slot chain of
  Section~\ref{sec:microbench} overstates that whole-queue gain. The
  baseline remainder times out at 300\,s for \(K=4,6,8\) in the recorded
  sweep, and at 120\,s in Table~\ref{tab:catalog}.
  \item Opus: the inner-product kernel.
  \item OpenCV: luma pixels.
  \item CMSIS-DSP: one \(\texttt{filter()}\) call. Extra trips from
  rounding 5 to 8 are slower, so this row emits \(P_1\).
  Table~\ref{tab:catalog} used a contract with no \texttt{bench.c}, so
  that row emitted the admitted \(\texttt{BLOCK}=8\) rewrite without a
  timing comparison (\S\ref{sec:implementation}).
\end{itemize}

\subsection{Misses of This Instance}
\label{sec:failures}

This subsection names where the three answers fail on \textsc{IntC}.
Table~\ref{tab:failures} is those misses.

\begin{table}
\caption{Misses of \textsc{IntC}: never generated, budget, slower
admit, weak harness, bound too small, driver, and unnamed cost.}
\label{tab:failures}
\begin{center}
\scriptsize
\begin{tabular}{llp{0.42\columnwidth}}
\toprule
Kind & Instance & What happens \\
\midrule
Never generated & linked queue, CRS-to-ELLPACK, swap-and-pop & not in the catalog \\
Budget & 1024-slot mask over 8016\,B & dropped, emit \(P_1\) \\
Slower admit & FIR \(\texttt{BLOCK}\) 5\(\to\)8 & extra trips, emit \(P_1\) \\
Weak harness & occupancy-only LIFO & FIFO harness refutes it \\
Bound too small & \(\texttt{\& 999}\) & proved at \(K=4\), refuted at \(K=10\) \\
Driver & slot chain vs ring \texttt{bench.c} & \(6.8\times\) vs \(2.3\times\), same rewrite \\
Unnamed cost & false sharing & not measured \\
\bottomrule
\end{tabular}
\end{center}
\end{table}

\textbf{Kind.} The cell is how the miss happens.

\textbf{Instance.} The cell is the extract or rewrite.

\textbf{What happens.} The cell is the outcome. In table order:
\begin{itemize}
  \item Never generated: a linked queue, a CRS-to-ELLPACK layout, and a
  swap-and-pop compaction are not catalog sites, so they never become
  candidates. A CRS-to-ELLPACK rewrite that kept the stored values would
  still stream the extra zeros ELLPACK pads.
  \item Budget: the 1024-slot mask proves and is 8208 bytes. The 8016-byte
  ceiling drops it and the emit is \(P_1\)
  (Section~\ref{sec:budget_eval}).
  \item Slower admit: FIR \(\texttt{BLOCK}\) 8 does extra trips on one
  \(\texttt{filter()}\) call, so ranking emits \(P_1\)
  (Section~\ref{sec:ring_eval}).
  \item Weak harness: an occupancy-only \(\mathit{Post}\) proves a LIFO
  stack that the FIFO harness refutes
  (Section~\ref{sec:microbench}).
  \item Bound too small: \(\texttt{\& 999}\) proves at \(K=4\) and is
  refuted at \(K=10\) (Section~\ref{sec:bounded_gap}).
  \item Driver: the isolated slot chain is \(6.8\times\). The ring
  \texttt{bench.c} is \(2.3\times\) on the same rewrite
  (Section~\ref{sec:microbench}).
  \item Unnamed cost: concurrent false sharing is not in any named
  driver, so it cannot reject a candidate.
\end{itemize}
